% Part: sets-functions-relations
% Chapter: sets
% Section: pairs-and-products

\documentclass[../../../include/open-logic-section]{subfiles}

\begin{document}

\olfileid{sfr}{set}{pai}
\olsection{যুগল, টিউপল ও কার্তেসীয় গুণফল}

\begin{explain}
সদস্যভিত্তিক সমতার নীতি থেকে বোঝা যায় যে, সেটের উপাদানগুলির কোনও
ক্রম নেই। তাই ক্রম প্রকাশ করতে চাইলে \emph{ক্রমযুগল} $\tuple{x, y}$
ব্যবহার করি। ক্রমহীন যুগল $\{x, y\}$-এ ক্রম বিবেচ্য নয়:
$\{x, y\} = \{y, x\}$। কিন্তু ক্রমযুগলে তা বিবেচ্য:
$x \neq y$ হলে $\tuple{x, y} \neq \tuple{y, x}$।

সেটতত্ত্বে ক্রমযুগলকে কীভাবে দেখব? মূল কথা হল, এই ধারণাটি অক্ষুণ্ণ
রাখতে চাই যে, দুটি ক্রমযুগল অভিন্ন হবে যদি এবং কেবল যদি তাদের প্রথম
উপাদান একই হয় এবং দ্বিতীয় উপাদানও একই হয়। অর্থাৎ:
\[
  \tuple{a, b}= \tuple{c, d}\text{ হবে যদি এবং কেবল যদি }a = c \text{ এবং }b=d.
\]
ভিনার--কুরাতোভস্কির সংজ্ঞার সাহায্যে সেটতত্ত্বে ক্রমযুগল সংজ্ঞায়িত করতে পারি।
\end{explain}

\begin{defn}[ক্রমযুগল]\ollabel{wienerkuratowski}
	$\tuple{a, b} = \{\{a\}, \{a, b\}\}$।
\end{defn}

\begin{prob}
	\olref[sfr][set][pai]{wienerkuratowski} ব্যবহার করে প্রমাণ করো যে,
	$\tuple{a, b}= \tuple{c, d}$ হবে যদি এবং কেবল যদি $a = c$ এবং $b=d$, উভয়ই হয়।
\end{prob}

\begin{explain}
ক্রমযুগলের সংজ্ঞা স্থির করার পরে তার সাহায্যে আরও সেট সংজ্ঞায়িত
করতে পারি। যেমন, কখনও দুটির বেশি বস্তুর ক্রমিত অনুক্রম দরকার হয়:
\emph{ত্রয়ী} $\tuple{x, y, z}$, \emph{চতুষ্টয়} $\tuple{x, y, z, u}$,
ইত্যাদি। ত্রয়ীকে বিশেষ ধরনের ক্রমযুগল হিসেবে ভাবতে পারি, যার প্রথম
উপাদান নিজেই ক্রমযুগল: $\tuple{x, y, z}$ হল $\tuple{\tuple{x, y},z}$।
চতুষ্টয়ের ক্ষেত্রেও তাই: $\tuple{x,y,z,u}$ হল
$\tuple{\tuple{\tuple{x,y},z},u}$; একইভাবে এগোনো যায়। সাধারণভাবে
আমরা \emph{ক্রমিত $n$-টিউপল} $\tuple{x_1, \dots, x_n}$-এর কথা বলি।

ক্রমযুগল কিংবা অন্য ক্রমিত $n$-টিউপলের কিছু সেট আমাদের কাজে লাগবে।
\end{explain}

\begin{defn}[কার্তেসীয় গুণফল]
প্রদত্ত সেট $A$ ও $B$-এর \emph{কার্তেসীয় গুণফল} $A \times B$-এর
সংজ্ঞা হল:
\[
  A \times B = \Setabs{\tuple{x, y}}{x \in A \text{ এবং } y \in B}.
\]
\end{defn}

\begin{ex}
$A = \{0, 1\}$ এবং $B = \{1, a, b\}$ হলে, তাদের গুণফল হল
\[
A \times B = \{ \tuple{0, 1}, \tuple{0, a}, \tuple{0, b},
    \tuple{1, 1}, \tuple{1, a}, \tuple{1, b} \}.
\]
\end{ex}

\begin{ex}
$A$ সেট হলে, $A$-এর নিজের সঙ্গে গুণফল $A \times A$-কে $A^2$-ও
লেখা হয়। এটি $x, y \in A$ এমন \emph{সমস্ত} যুগল $\tuple{x, y}$-এর
সেট। সমস্ত ত্রয়ী $\tuple{x, y, z}$-এর সেট হল $A^3$, এবং এভাবেই
চলতে থাকে। আমরা একটি পুনরাবৃত্তিমূলক সংজ্ঞা দিতে পারি:
\begin{align*}
  A^1 & = A\\
  A^{k+1} & = A^k \times A
\end{align*}
\end{ex}

\begin{prob}
$\{1, 2, 3\}^3$-এর সমস্ত !!{element}sের তালিকা তৈরি করো।
\end{prob}

\begin{prop}\ollabel{cardnmprod}
$A$-এর $n$টি !!{element}s এবং $B$-এর $m$টি !!{element}s থাকলে,
$A \times B$-এর $n\cdot m$টি উপাদান থাকে।
\end{prop}

\begin{proof}
$A$-এর প্রতিটি !!{element}~$x$-এর জন্য $\tuple{x, y} \in A \times B$
আকারের $m$টি !!{element}s আছে। ধরি,
$B_x = \Setabs{\tuple{x, y}}{y \in B}$। যখনই $x_1 \neq x_2$,
তখনই $\tuple{x_1, y} \neq \tuple{x_2, y}$;
তাই $B_{x_1} \cap B_{x_2} = \emptyset$। কিন্তু $A = \{x_1, \dots, x_n\}$
হলে $A \times B = B_{x_1} \cup \dots \cup B_{x_n}$।
কাজেই এতে $n\cdot m$টি !!{element}s আছে।

এটি দৃশ্যায়িত করতে $A \times B$-এর !!{element}s-কে ছকের আকারে সাজাও:
\[
\begin{array}{rcccc}
  B_{x_1} = & \{\tuple{x_1, y_1} & \tuple{x_1, y_2} & \dots & \tuple{x_1, y_m}\}\\
  B_{x_2} = & \{\tuple{x_2, y_1} & \tuple{x_2, y_2} & \dots & \tuple{x_2, y_m}\}\\
  \vdots & & \vdots\\
  B_{x_n} = & \{\tuple{x_n, y_1} & \tuple{x_n, y_2} & \dots & \tuple{x_n, y_m}\}
\end{array}
\]
$x_i$-গুলি সব আলাদা এবং $y_j$-গুলিও সব আলাদা। তাই এই ছকের
কোনও দুটি যুগল এক নয়, আর মোট যুগলের সংখ্যা $n\cdot m$।
\end{proof}

\begin{prob}
$k$-এর উপর আরোহ পদ্ধতি ব্যবহার করে দেখাও যে, সমস্ত $k \ge 1$-এর জন্য,
$A$-এর $n$টি !!{element}s থাকলে $A^k$-এর $n^k$টি !!{element}s থাকে।
\end{prob}

\begin{ex}
$A$ সেট হলে, $A$-এর !!{element}sের যে-কোনও অনুক্রমকে $A$-এর উপর
একটি \emph{শব্দ} বলে। অনুক্রমকে $A$-এর !!{element}sের একটি
$n$-টিউপল হিসেবে ভাবা যায়। যেমন, $A = \{a, b, c\}$ হলে ``$bac$''
অনুক্রমটিকে ত্রয়ী~$\tuple{b, a, c}$ হিসেবে ভাবতে পারি। শব্দ, অর্থাৎ
প্রতীকের অনুক্রম, কম্পিউটার বিজ্ঞানে অত্যন্ত গুরুত্বপূর্ণ। প্রচলিত রীতি
অনুযায়ী $A$-এর !!{element}s-কে দৈর্ঘ্য~$1$-এর অনুক্রম এবং
$\emptyset$-কে দৈর্ঘ্য~$0$-এর অনুক্রম হিসেবে গণ্য করি।
তাহলে $A$-এর উপর \emph{সমস্ত} শব্দের সেট হল
\[
A^* = \{\emptyset\} \cup A \cup A^2 \cup A^3 \cup \dots
\]
\end{ex}

\end{document}
